% !PRINTED 570
\weberpage{570}
%<<CONT
und da $\sin\frac{11\pi}{34}$ grösser ist als $\sin\frac{3\pi}{34}$, so ist die Summe dieser
beiden Ausdrücke und damit die ganze Summe $(-1,\eta)$ positiv % [sic] the printed authority ends the sentence without punctuation

Es ergiebt sich ferner für den reellen Theil von $(i,\eta)$

\[
\eta - \eta_2 + \eta_4 - \eta_6 = 2\cos\frac{2\pi}{17} + 2\cos\frac{\pi}{17} + 2\cos\frac{8\pi}{17} - 2\cos\frac{4\pi}{17},
\]

was unmittelbar als positiv erkannt wird; endlich für den reellen
Theil von $(\alpha,\eta)$

\[
\eta - \eta_4 + \frac{1}{\sqrt{2}}(\eta_1 - \eta_3 - \eta_5 + \eta_7)
= 2\cos\frac{2\pi}{17} - 2\cos\frac{8\pi}{17} + \sqrt{2}\left(\cos\frac{6\pi}{17} + \cos\frac{3\pi}{17} + \cos\frac{7\pi}{17} - \cos\frac{5\pi}{17}\right),
\]

was sich gleichfalls als positiv erweist.

\begin{center}
\S.\,170.

\gesperrt{Eigenschaften der Zahlen} $\psi$.
\end{center}

Die im vorigen Paragraphen abgeleiteten Formeln für die
Resolventen gestatten eine Fülle von Anwendungen, auch ohne
eine genaue Bestimmung der Primzahl $n$, für specielle Annahmen
über die Zahl $e$; sie geben also zahlentheoretische Sätze für
ganze Kategorien von Primzahlen.

Für die Algebra liefern sie uns die Lösungen cyklischer
Gleichungen von vorgeschriebenem Grade in beliebiger Menge.

Der Ableitung der für die Algebra wichtigsten dieser Anwendungen
schicken wir einige allgemeine Betrachtungen über
die Functionen $\psi$ voraus, die wir im vorigen Paragraphen definirt
haben.

Zur Berechnung der Zahlen $\psi_{\lambda,\mu}(\varepsilon)$ [\S.\,169, (5)] bedient man
sich, wie wir im Falle $n = 17$ gesehen haben, einer Indextabelle.
Es bestehen aber zwischen diesen Functionen Relationen, durch
die sie sich auf eine geringere Zahl zurückführen lassen, worüber
\gesperrt{Jacobi} merkwürdige Untersuchungen angestellt hat\footnote{Vergl. \gesperrt{Kronecker}, Zur Theorie der \gesperrt{Abel'}schen Gleichungen. Journ. f. Mathem. Bd. 93.}.

Die Function $\psi_{\lambda,\mu}(\varepsilon)$ ändert sich nach ihrer Definition
nicht, wenn $\lambda,\mu$ vertauscht werden, obwohl ihr Ausdruck
%CONT>>
